\uuid{tV8a}
\titre{Limite en un point}
\auteur{bodin ; hanani ; mzari}
\organisation{exo7}
\licence{CC BY-NC-SA 4.0}
\datecreate{2020-01-20}
\niveau{L1}
\module{Analyse}
\chapitre{Continuité, limite et étude de fonctions réelles}
\sousChapitre{Limite de fonctions}
\difficulte{2}
\theme{Limites des fonctions réelles, limite en un point}
\exo7qcmid{lille-1.limites.019}

\contenu{
\begin{question}
\qclassification{123.03}
\qtags{motcle=limite en un point}
Soit $f(x)= \frac{\sin(3x)}{\sin(4x)}$. Quelles sont les assertions vraies ?
\begin{answers}

    \bad{$f$ n'admet pas de limite en $0$}

    \good{$\lim_{x\to 0} f(x)=\frac{3}{4}$}

    \bad{$\lim_{x\to 0} f(x)= \frac{4}{3}$}

    \bad{$\lim_{x\to 0} f(x)=0$}

\end{answers}
\begin{explanations}
$\lim_{x\to 0}\frac{\sin x}{x} = 1,$ donc $\lim_{x\to 0} f(x)= \lim_{x\to 0}\frac{3x}{4x} \cdot  \frac{\sin (3x)}{3x}\cdot  \frac{4x}{\sin(4x)} = \frac{3}{4}$.
\end{explanations}
\end{question}
}
